\sect{कार्त्तिक-शुक्ल-प्रबोधिनी-एकादशी-माहात्म्यम्}
\label{sec:padma-karttika-shukla-prabodhini}

\ganapatyadiDhyanam
\padmaPuranam

\dnsub{अथ त्रिषष्टितमोऽध्यायः॥६३}

\uvacha{युधिष्ठिर उवाच}

\twolineshloka
{श्रुतं रमाया माहात्म्यं त्वत्तः कृष्ण यथातथम्}
{कार्तिके शुक्लपक्षे या तां मे कथय मानद}% ॥१॥

\uvacha{श्रीकृष्ण उवाच}

\twolineshloka
{शृणु राजन् प्रवक्ष्यामि शुक्ले चोर्जदले तु या}
{सा यथा नारदे प्रोक्ता ब्रह्मणा लोककारिणा}% ॥२॥

\uvacha{नारद उवाच}

\twolineshloka
{प्रबोधिन्याश्च माहात्म्यं वद विस्तरतो मम}
{यस्यां जागर्ति गोविन्दो धर्मकर्मप्रवर्तकः}% ॥३॥

\uvacha{ब्रह्मोवाच}

\twolineshloka
{प्रबोधिन्याश्च माहात्म्यं पापघ्नं पुण्यवर्धनम्}
{मुक्तिप्रदं सुबुद्धीनां शृणुष्व मुनिसत्तम}% ॥४॥

\twolineshloka
{तावद्गर्जन्ति तीर्थानि आसमुद्रसरांसि च}
{यावत् प्रबोधिनी विष्णोस्तिथिर्नाऽऽयाति कार्तिके}% ॥५॥

\twolineshloka
{तावद्गर्जन्ति विप्रेन्द्र गङ्गाभागीरथी क्षितौ}
{यावन्नायाति पापघ्नी कार्तिके हरिबोधिनी}% ॥६॥

\twolineshloka
{अश्वमेधसहस्राणि राजसूयशतानि च}
{एकेनैवोपवासेन प्रबोधिन्या लभेन्नरः}% ॥७॥

\twolineshloka
{यद् दुर्लभं यदप्राप्यं त्रैलोक्यस्य न गोचरम्}
{तदप्यप्रार्थितं पुत्रं ददाति हरिबोधिनी}% ॥८॥

\twolineshloka
{ऐश्वर्यं सम्पदं प्रज्ञां राज्यं च सुखसम्पदः}
{ददात्युपोषिता भक्त्या जनेभ्यो हरिबोधिनी}% ॥९॥

\twolineshloka
{मेरुमन्दरमात्राणि पापान्युक्तानि यानि च}
{एकेनैवोपवासेन दहते पापनाशिनी}% ॥१०॥

\twolineshloka
{पूर्वजन्मसहस्रेषु यत् पापं समुपार्जितम्}
{निशि जागरणं चास्या दहते तूलराशिवत्}% ॥११॥

\twolineshloka
{उपवासं प्रबोधिन्यां यः करोति स्वभावतः}
{विधिवन्मुनिशार्दूल यथोक्तं लभते फलम्}% ॥१२॥

\twolineshloka
{यथोक्तं कुरुते यस्तु विधिवत्सुकृतं नरः}
{स्वल्पं मुनिवरश्रेष्ठ मेरुतुल्यं भवेत्फलम्}% ॥१३॥

\twolineshloka
{विधिहीनं तु यः कुर्यात्सुकृतं मेरुमात्रकम्}
{अणुमात्रं तदाप्नोति फलं धर्मस्य नारद}% ॥१४॥

\threelineshloka
{ये ध्यायन्ति मनोवृत्या ये करिष्यन्ति बोधिनीम्}
{वसन्ति पितरो हृष्टा विष्णुलोके च तस्य वै}% ॥१५॥
{विमुक्ता नारकैर्दुःखैर्याति विष्णोः परं पदम्}

\twolineshloka
{कृत्वा तु पातकं घोरं ब्रह्महत्यादिकं नरः}% ॥१६॥
{कृत्वा तु जागरं विष्णोर्धौतपापो भवेन्नरः}

\twolineshloka
{दुष्प्राप्यं यत्फलं विप्र अश्वमेधादिकैर्मखैः}% ॥१७॥
{प्राप्यते तत् सुखेनैव प्रबोधिन्यास्तु जागरे}

\twolineshloka
{आप्लुत्य सर्वतीर्थेषु प्रदत्त्वा काञ्चनं महीम्}% ॥१८॥
{तत्फलं समवाप्नोति यत्कृत्वा जागरं हरेः}

\twolineshloka
{जातः स एव सुकृती कुलं तेनैव पावितम्}% ॥१९॥
{कार्तिके मुनिशार्दूल कृता येन प्रबोधिनी}

\twolineshloka
{यथा ध्रुवं नृणां मृत्युर्धनं गात्रं तथाऽध्रुवम्}% ॥२०॥
{इति ज्ञात्वा मुनिश्रेष्ठ कर्तव्यं वैष्णवं दिनम्}

\twolineshloka
{यानि कानि च तीर्थानि त्रैलोक्ये सम्भवन्ति च}% ॥२१॥
{तानि तस्य गृहे सम्यग्यः करोति प्रबोधिनीम्}

\twolineshloka
{किं तस्य बहुभिः पुण्यैः कृता येन प्रबोधिनी}% ॥२२॥
{पुत्रपौत्रप्रदा ह्येषा कार्तिके हरिबोधिनी}

\twolineshloka
{स ज्ञानी च स योगी च स तपस्वी जितेन्द्रियः}% ॥२३॥
{भोगो मोक्षश्च तस्यास्ति उपास्ते हरिबोधिनीम्}

\twolineshloka
{विष्णोः प्रियतरा ह्येषा धर्मसारसहायिनी}% ॥२४॥
{यः करोति नरो भक्त्या भुक्तिभाक्स भवेन्नरः}

\twolineshloka
{प्रबोधिनीमुपोषित्वा गर्भे न विशते नरः}% ॥२५॥
{सर्वधर्मान् परित्यज्य तस्मात् कुर्वीत नारद}

\twolineshloka
{कर्मणा मनसा वाचा पापं यत्समुपार्जितम्}% ॥२६॥
{तत्क्षालयति गोविन्दः प्रबोधिन्यां तु जागरे}

\twolineshloka
{स्नानं दानं जपः पूजां समुद्दिश्य जनार्दनम्}% ॥२७॥
{नरो यत्कुरुते वत्स प्रबोधिन्यां तदक्षयम्}

\twolineshloka
{येऽर्चयन्ति नरास्तस्यां भक्त्या देवं च माधवम्}% ॥२८॥
{समुपोष्य प्रमुच्यन्ते पापैस्तैः शतजन्मजैः}

\twolineshloka
{महाव्रतमिदं पुत्र महापापौघनाशनम्}% ॥२९॥
{प्रबोधवासरं विष्णोर्विधिवत्समुपोषयेत्}

\twolineshloka
{व्रतेनानेन देवेशं परितोष्य जनार्दनम्}% ॥३०॥
{विराजयन्दिशः सर्वाः प्रयाति हरिमन्दिरम्}

\twolineshloka
{कर्तव्यैषा प्रयत्नेन नरैः कान्तिधनार्थिभिः}% ॥३१॥
{बाल्ये यत्सञ्चितं पापं यौवने वार्धके तथा}

\twolineshloka
{शतजन्मकृतं पापं स्वल्पं वा यदि वा बहु}% ॥३२॥
{तत्क्षालयति गोविन्दश्चास्यामभ्यर्चितो नृणाम्}

\twolineshloka
{धनधान्यवहा पुण्या सर्वपापहरा परा}% ॥३३॥
{तामुपोष्य हरेर्भक्त्या दुर्लभं न भवेत्क्वचित्}

\twolineshloka
{चन्द्रसूर्योपरागे च यत्फलं परिकीर्तितम्}% ॥३४॥
{तत् सहस्रगुणं प्रोक्तं प्रबोधिन्यां प्रजागरे}

\twolineshloka
{स्नानं दानं जपो होमः स्वाध्यायोऽभ्यर्चनं हरेः}% ॥३५॥
{तत् सर्वं कोटिगुणितं प्रबोधिन्यां कृतं तु यत्}

\twolineshloka
{जन्मप्रभृति यत् पुण्यं नरेणोपार्जितं भवेत्}% ॥३६॥
{वृथा भवति तत् सर्वमकृत्वा कार्तिके व्रतम्}

\twolineshloka
{अकृत्वा नियमं विष्णोः कार्तिकं यः क्षिपेन्नरः}% ॥३७॥
{न जन्मार्जितपुण्यस्य फलं प्राप्नोति नारद}

\twolineshloka
{तस्मात् सर्वप्रयत्नेन देवदेवं जनार्दनम्}% ॥३८॥
{उपसेवेत विप्रेन्द्र सर्वकामफलप्रदम्}

\twolineshloka
{परान्नं वर्जयेद्यस्तु कार्तिके विष्णुतत्परः}% ॥३९॥
{परान्नवर्जनाद्वत्स चान्द्रायणफलं लभेत्}

\twolineshloka
{नित्यं शास्त्रविनोदेन कार्तिके मधुसूदनः}% ॥४०॥
{स दहेत् सर्वपापानि यज्ञायुतफलं लभेत्}

\twolineshloka
{न तथा तुष्यते यज्ञैर्न दानैर्वाजपादिभिः}% ॥४१॥
{यथा शास्त्रकथालापैः कार्तिके मधुसूदनः}

\twolineshloka
{ये कुर्वन्ति कथां विष्णोर्ये शृण्वन्ति शुभान्विताः}% ॥४२॥
{श्लोकार्धं श्लोकपादं वा कार्तिके गोशतं फलम्}

\twolineshloka
{सर्वधर्मान् परित्यज्य कार्तिके केशवाग्रतः}% ॥४३॥
{शास्त्रावधारणं कार्यं श्रोतव्यं च महामुने}

\twolineshloka
{श्रेयसा लोभबुद्ध्या च यः करोति हरेः कथाम्}% ॥४४॥
{कार्तिके मुनिशार्दूल कुलानां तारयेच्छतम्}

\twolineshloka
{नियमेन नरो यस्तु शृणुते वैष्णवीं कथाम्}% ॥४५॥
{कार्तिके तु विशेषेण गोसहस्रफलं लभेत्}

\twolineshloka
{प्रबोधवासरे विष्णोः शृणुते यो हरेः कथाम्}% ॥४६॥
{सप्तद्वीपवती दाने तत्फलं लभते मुने}

\twolineshloka
{श्रुत्वा विष्णुकथां दिव्यां येऽर्चयन्ति कथाविदम्}% ॥४७॥
{स्वशक्त्या मुनिशार्दूल तेषां लोकोऽक्षयः स्मृतः}

\twolineshloka
{गीतशास्त्रविनोदेन कार्तिकं यो नयेन्नरः}% ॥४८॥
{न तस्य पुनरावृत्तिर्मया दृष्टा कलिप्रिय}

\twolineshloka
{गीतं नृत्यं च वाद्यं च भव्यां विष्णुकथां मुने}% ॥४९॥
{यः करोति स पुण्यात्मा त्रैलोक्योपरि संस्थितः}

\twolineshloka
{बहुपुष्पैर्बहुफलैः कर्पूरागुरुकुङ्कुमैः}% ॥५०॥
{हरेः पूजा विधातव्या कार्तिके बोधवासरे}

\twolineshloka
{यस्मात् पुण्यमसङ्ख्यातं प्राप्यते मुनिसत्तम}% ॥५१॥
{फलैर्नानाविधैर्द्रव्यैः प्रबोधिन्यां तु जागरे}

\onelineshloka
{शङ्खे तोयं समादाय अर्घो देयो जनार्दने}% ॥५२॥

\twolineshloka
{यत्फलं सर्वतीर्थेषु सर्वदानेषु यत्फलम्}
{तत्फलं कोटिगुणितं दत्त्वाऽर्घं बोधवासरे}% ॥५३॥

\twolineshloka
{गुरुपूजा ततः कार्या भोजनाच्छादनादिभिः}
{दक्षिणाभिश्च देवर्षे तुष्ट्यर्थं चक्रपाणिनः}% ॥५४॥

\twolineshloka
{भागवतं शृणुते यस्तु पुराणं च पठेन्नरः}
{प्रत्यक्षरं भवेत् तस्य कपिलादानजं फलम्}% ॥५५॥

\twolineshloka
{कार्त्तिके मुनिशार्दूल स्वशक्त्या वैष्णवं व्रतम्}
{यः करोति यथोक्तं तु मुक्तिस्तस्य सुनिश्चला}% ॥५६॥

\twolineshloka
{केतक्या एकपत्रेण पूजितो गरुडध्वजः}
{समाः सहस्रं सुप्रीतो भवति मधुसूदनः}% ॥५७॥

\twolineshloka
{अगस्तिकुसुमैर्देवं पूजयेद्यो जनार्दनम्}
{दर्शनात् तस्य देवर्षे नरकाग्निः प्रणश्यति}% ॥५८॥

\twolineshloka
{मुनिपुष्पार्चितो विष्णुः कार्तिके पुरुषोत्तमः}
{ददात्यभिमतान्कामान्शशिसूर्यग्रहे यथा}% ॥५९॥

\twolineshloka
{विहाय सर्वपुष्पाणि मुनिपुष्पेण केशवम्}
{कार्तिके योऽर्चयेद् भक्त्या वाजिमेधफलं लभेत्}% ॥६०॥

\twolineshloka
{तुलसीदलानि पुष्पाणि ये यच्छन्ति जनार्दने}
{कार्तिके सकलं वत्स पापं जन्मायुतं दहेत्}% ॥६१॥

\twolineshloka
{दृष्टा स्पृष्टाऽथ वा ध्याता कीर्तिता नामतस्तु ता}
{रोपिता सिञ्चिता नित्यं पूजिता तुलसी शुभा}% ॥६२॥

\twolineshloka
{नवधा तुलसीभक्तिं ये कुर्वन्ति दिनेदिने}
{युगकोटिसहस्राणि तन्वन्ति सुकृतं मुने}% ॥६३॥

\twolineshloka
{यावच्छाखाप्रशाखाभिर्बीजपुष्पदलैर्मुने}
{रोपिता तुलसी पुम्भिर्वर्धते वसुधातले}% ॥६४॥

\twolineshloka
{तेषां वंशे तु ये जाता ये भविष्यन्ति ये गताः}
{आकल्पवर्षसाहस्रं तेषां वासो हरेर्गृहे}% ॥६५॥

\twolineshloka
{यत्फलं सर्वपुष्पेषु सर्वपत्रेषु नारद}
{तुलसीदलेन चैकेन कार्तिके प्राप्यते तु तत्}% ॥६६॥

\twolineshloka
{सम्प्राप्तं कार्त्तिकं दृष्ट्वा नियमेन जनार्दनः}
{पूजनीयो महाविष्णुः कोमलैस्तुलसीदलैः}% ॥६७॥

\twolineshloka
{इष्ट्वा क्रतुशतैर्देवान्दत्त्वा दानान्यनेकशः}
{तुलसीदलैस्तु तत् पुण्यं कार्तिके केशवार्चने}% ॥६८॥

॥इति श्रीपाद्मे महापुराणे पञ्चपञ्चाशत्साहस्र्यां संहितायामुत्तरखण्डे उमापतिनारदसंवादान्तर्गतकृष्णयुधिष्ठिरसंवादे कार्त्तिक-शुक्ल-प्रबोधिनी-एकादशी-माहात्म्यं नाम त्रिषष्टितमोऽध्यायः॥६३॥

\genMangalaShloka

\hyperref[sec:ekadashi_mahatmyam_padma_puranam]{\closesub}
\clearpage

